\documentclass[a4paper,12pt]{article}

\usepackage[T2A]{fontenc}
\usepackage[utf8]{inputenc}
\usepackage[russian]{babel}
\usepackage{amsmath,amssymb,mathtools}
\usepackage{helvet}
\renewcommand{\familydefault}{\sfdefault}
\usepackage{icomma}
\usepackage[a4paper,margin=2.2cm]{geometry}
\usepackage{microtype}
\usepackage{tikz}
\usepackage{tkz-euclide}
\usepackage{tabularx,booktabs}
\usepackage{enumitem}
\usepackage{hyperref}
\usepackage{parskip}
\usepackage{fancyhdr}
\usepackage{setspace}
\usepackage{xcolor}

\definecolor{accent}{HTML}{008080}
\definecolor{darkaccent}{HTML}{004D4D}
\definecolor{textgray}{HTML}{333333}
\definecolor{softgray}{HTML}{E9EEEE}

\hypersetup{
  colorlinks=true,
  linkcolor=darkaccent,
  urlcolor=darkaccent,
  pdftitle={Чек-лист по планиметрии},
  pdfauthor={Лёвин Артём Александрович}
}

\onehalfspacing
\setlength{\parindent}{0pt}
\setlength{\headheight}{25pt}

\setlist[itemize]{
  leftmargin=1.6em,
  itemsep=0.25em,
  topsep=0.35em
}

\setlist[enumerate]{
  leftmargin=2em,
  itemsep=0.45em,
  topsep=0.45em
}

\pagestyle{fancy}
\fancyhf{}
\fancyhead[L]{\small\textcolor{textgray}{Планиметрия: треугольник и окружность}}
\fancyhead[R]{\small\textcolor{textgray}{Анна Трапезникова}}
\fancyfoot[C]{\small\thepage}
\renewcommand{\headrulewidth}{0.3pt}
\renewcommand{\headrule}{
  \hbox to\headwidth{
    \color{softgray}
    \leaders\hrule height \headrulewidth\hfill
  }
}

\newcommand{\sectionrule}{
  \par
  \vspace{0.2em}
  \noindent
  \textcolor{accent}{\rule{\linewidth}{0.7pt}}
  \vspace{0.35em}
}

\newcommand{\key}[1]{\textbf{\textcolor{darkaccent}{#1}}}
\newcommand{\formula}[1]{\[\boxed{#1}\]}

\begin{document}

% ============================================================
% ТИТУЛЬНЫЙ ЛИСТ
% ============================================================

\begin{titlepage}
\thispagestyle{empty}
\centering

\vspace*{2.5cm}

{\Large\textcolor{darkaccent}{\textbf{ЧЕК-ЛИСТ}}\par}

\vspace{0.8cm}

{\huge\bfseries Планиметрия:\par}

\vspace{0.2cm}

{\LARGE\bfseries
прямоугольный треугольник, окружность, хорды\\
и теорема Птолемея
\par}

\vspace{1.2cm}

{\large ЕГЭ, профильная математика\par}

\vfill

{\large Лёвин Артём Александрович\par}

\vspace{0.15cm}

{\large эксклюзивно для Анны Трапезниковой\par}

\vspace{0.6cm}

{\large \today\par}

\vspace{1.6cm}

\begin{tikzpicture}[x=1cm,y=1cm]
  \draw[
    accent,
    line width=1.1pt
  ]
    (-6,0)
    .. controls (-3.5,0.9) and (-1.5,-0.8) .. (0,0)
    .. controls (1.7,0.9) and (3.7,-0.6) .. (6,0.15);
\end{tikzpicture}

\vspace*{0.8cm}

\end{titlepage}

% ============================================================
% ОСНОВНОЙ ЧЕК-ЛИСТ
% ============================================================

\section*{1. Как читать и оформлять геометрическую задачу}
\sectionrule

\begin{itemize}
  \item Сначала выпишите \key{дано} и \key{что требуется найти или доказать}.
  Это дисциплинирует перевод условия с русского языка на математический.

  \item На чертеже подписывайте точки, известные длины и углы.
  Не приписывайте фигуре свойства, которых нет в условии.

  \item Если сказано только «пятиугольник вписан в окружность»,
  нельзя автоматически считать его правильным.

  \item Запись вида «\(\sin\)» без указания угла некорректна.
  Пишите, например,
  \[
  \sin\angle BAC.
  \]

  \item Чертёж полезно делать приблизительно согласованным
  с отношениями длин: более длинная сторона должна выглядеть
  более длинной.
\end{itemize}

\section*{2. Прямоугольный треугольник: базовый набор}
\sectionrule

Пусть \(\triangle ABC\) прямоугольный,
\[
\angle ACB=90^\circ.
\]
Тогда \(AB\) --- гипотенуза, \(AC\) и \(BC\) --- катеты.

\formula{AB^2=AC^2+BC^2}

Для острого угла \(\angle BAC\):

\formula{
\sin\angle BAC=\frac{BC}{AB},
\qquad
\cos\angle BAC=\frac{AC}{AB},
\qquad
\tg\angle BAC=\frac{BC}{AC}
}

\key{Полезная связь для дополнительных углов:}
\[
\sin\alpha=\cos(90^\circ-\alpha),
\qquad
\cos\alpha=\sin(90^\circ-\alpha).
\]

\subsection*{Пример 1. Удобные вычисления вместо громоздкой арифметики}

В прямоугольном \(\triangle ABC\):
\[
\angle ACB=90^\circ,
\qquad
AC=5\sqrt{21},
\qquad
BC=10.
\]
Найдём \(\sin\angle BAC\).

\begin{center}
\begin{tikzpicture}[scale=0.9]

  \tkzDefPoint(0,0){C}
  \tkzDefPoint(5.2,0){A}
  \tkzDefPoint(0,2.1){B}

  \tkzDrawPolygon[
    line width=1pt,
    color=accent
  ](A,B,C)

  \tkzDrawPoints(A,B,C)

  \tkzLabelPoints[below](A,C)
  \tkzLabelPoints[above left](B)

  \tkzMarkRightAngle[size=0.35](A,C,B)

  \tkzLabelSegment[below](C,A){$5\sqrt{21}$}
  \tkzLabelSegment[left](C,B){$10$}

\end{tikzpicture}
\end{center}

По теореме Пифагора:
\[
AB^2=(5\sqrt{21})^2+10^2.
\]

Вместо вычисления \(525+100\) удобно сгруппировать:
\[
AB^2
=
25\cdot21+25\cdot4
=
25(21+4)
=
25^2.
\]

Следовательно,
\[
AB=25.
\]

По определению синуса:
\[
\sin\angle BAC
=
\frac{BC}{AB}
=
\frac{10}{25}
=
\frac25
=
0{,}4.
\]

\key{Практический приём.}
Перед прямым вычислением ищите общий множитель,
полный квадрат или удобную группировку.

\section*{3. Окружность: касательная, радиус, центральные углы и дуги}
\sectionrule

\begin{itemize}

  \item Радиус, проведённый в точку касания,
  перпендикулярен касательной:
  \[
  OA\perp AC
  \quad\Rightarrow\quad
  \angle OAC=90^\circ.
  \]

  \item Центральный угол равен градусной мере дуги,
  на которую он опирается:
  \[
  \angle AOD=\overset{\frown}{AD}.
  \]

  \item Если \(BD\) --- диаметр, то дуга полуокружности \(BD\)
  имеет величину
  \[
  180^\circ.
  \]

  \item Вписанный угол равен половине градусной меры дуги,
  на которую он опирается:
  \[
  \angle ABC
  =
  \frac12\overset{\frown}{AC}.
  \]

\end{itemize}

\subsection*{Пример 2. Касательная и секущая через центр}

Пусть \(CA\) касается окружности в точке \(A\).
Прямая \(CO\) проходит через центр \(O\)
и пересекает окружность в точках \(B\) и \(D\).

Дано:
\[
\overset{\frown}{AD}=116^\circ.
\]

Найти:
\[
\angle ACO.
\]

\begin{center}
\begin{tikzpicture}[scale=1]

  \tkzDefPoint(0,0){O}
  \tkzDefPoint(2,0){D}
  \tkzDefPoint(-2,0){B}
  \tkzDefPoint(-0.876,1.798){A}
  \tkzDefPoint(-4.565,0){C}

  \tkzDrawCircle[
    line width=1pt,
    color=accent
  ](O,D)

  \tkzDrawSegments[
    line width=1pt,
    color=accent
  ](C,A O,A C,D)

  \tkzDrawPoints(A,B,C,D,O)

  \tkzLabelPoints[above](A)
  \tkzLabelPoints[below](B,C,D,O)

  \tkzMarkRightAngle[size=0.3](C,A,O)

  \tkzMarkAngle[size=0.55](A,O,D)
  \tkzLabelAngle[pos=0.85](A,O,D){$116^\circ$}

\end{tikzpicture}
\end{center}

\key{Способ 1. Через базовые свойства.}

Центральный угол равен соответствующей дуге:
\[
\angle AOD=116^\circ.
\]

Точки \(C,O,D\) лежат на одной прямой, поэтому
углы \(\angle AOC\) и \(\angle AOD\) смежные:
\[
\angle AOC
=
180^\circ-116^\circ
=
64^\circ.
\]

Радиус \(OA\) перпендикулярен касательной \(AC\):
\[
\angle OAC=90^\circ.
\]

В \(\triangle AOC\):
\[
\angle ACO
=
180^\circ-90^\circ-64^\circ
=
26^\circ.
\]

\[
\boxed{\angle ACO=26^\circ}
\]

\key{Способ 2. Через формулу угла между касательной и секущей.}

Угол между касательной и секущей,
проведёнными из одной внешней точки,
равен половине разности заключённых между ними дуг:

\formula{
\angle ACD
=
\frac12
\left(
\overset{\frown}{AD}
-
\overset{\frown}{AB}
\right)
}

Поскольку \(BD\) --- диаметр,
\[
\overset{\frown}{AB}
=
180^\circ-116^\circ
=
64^\circ.
\]

Следовательно,
\[
\angle ACD
=
\frac12(116^\circ-64^\circ)
=
26^\circ.
\]

\key{Универсальный дополнительный ход.}
Если в задаче есть касательная,
проведите радиус в точку касания.
Так появляется прямой угол и часто возникает
полезный прямоугольный треугольник.

\section*{4. Равные хорды и равные дуги}
\sectionrule

В одной и той же окружности равные хорды стягивают
равные соответствующие дуги:

\formula{
AB=CD
\quad\Longleftrightarrow\quad
\overset{\frown}{AB}
=
\overset{\frown}{CD}
}

Отсюда следуют важные рабочие свойства:

\begin{itemize}

  \item равные дуги стягиваются равными хордами;

  \item вписанные углы, опирающиеся на одну и ту же дугу,
  равны;

  \item вписанные углы, опирающиеся на равные дуги,
  также равны.

\end{itemize}

Пусть пятиугольник \(ABCDE\) вписан в окружность и
\[
AB=CD,
\qquad
BC=DE.
\]

Тогда
\[
\overset{\frown}{AB}
=
\overset{\frown}{CD},
\qquad
\overset{\frown}{BC}
=
\overset{\frown}{DE}.
\]

Складывая равные дуги, получаем
\[
\overset{\frown}{ABC}
=
\overset{\frown}{CDE}.
\]

Эти дуги стягиваются хордами \(AC\) и \(CE\),
поэтому

\formula{AC=CE}

\section*{5. Вписанная трапеция}
\sectionrule

\key{Теорема.}
Если трапеция вписана в окружность,
то она равнобедренная.

Если
\[
AD\parallel BC
\]
и \(ABCD\) вписан в окружность, то
\[
AB=CD.
\]

У равнобедренной трапеции диагонали равны:
\[
AC=BD.
\]

Следовательно, для вписанной трапеции:
\[
AD\parallel BC
\quad\Rightarrow\quad
AB=CD,
\qquad
AC=BD.
\]

\subsection*{Как получить трапецию из равных хорд}

Пусть \(ABCD\) --- вписанный четырёхугольник и
\[
AB=CD.
\]

Тогда
\[
\overset{\frown}{AB}
=
\overset{\frown}{CD}.
\]

Следовательно, равны вписанные углы,
опирающиеся на эти дуги:
\[
\angle ADB
=
\angle DBC.
\]

Эти углы являются накрест лежащими
при секущей \(BD\), поэтому
\[
AD\parallel BC.
\]

Значит, \(ABCD\) --- трапеция.
Поскольку она вписана в окружность,
она равнобедренная.

\section*{6. Теорема Птолемея}
\sectionrule

Для любого четырёхугольника \(ABCD\),
вписанного в окружность,

\formula{
AC\cdot BD
=
AB\cdot CD
+
BC\cdot AD
}

То есть:

\[
\boxed{
\text{произведение диагоналей}
=
\text{сумма произведений противоположных сторон}
}
\]

\begin{center}
\begin{tikzpicture}[scale=0.9]

  \tkzDefPoint(0,0){O}
  \tkzDefPoint(2.8,0){R}

  \tkzDefPoint(2.35,1.52){A}
  \tkzDefPoint(-0.49,2.76){B}
  \tkzDefPoint(-2.56,-1.14){C}
  \tkzDefPoint(1.35,-2.45){D}

  \tkzDrawCircle[
    line width=1pt,
    color=accent
  ](O,R)

  \tkzDrawPolygon[
    line width=1pt,
    color=accent
  ](A,B,C,D)

  \tkzDrawSegments[
    line width=0.9pt,
    color=darkaccent
  ](A,C B,D)

  \tkzDrawPoints(A,B,C,D)

  \tkzLabelPoints[above right](A)
  \tkzLabelPoints[above left](B)
  \tkzLabelPoints[below left](C)
  \tkzLabelPoints[below right](D)

\end{tikzpicture}
\end{center}

\key{Сигнал к применению.}
Если перед Вами вписанный четырёхугольник,
известны несколько его сторон
и требуется диагональ или ещё одна сторона,
проверьте возможность применения теоремы Птолемея.

\subsection*{Пример 3. Двойное применение теоремы Птолемея}

Пятиугольник \(ABCDE\) вписан в окружность.

Дано:
\[
AB=CD=3,
\qquad
BC=DE=4,
\qquad
AD=6.
\]

Найти:
\[
BE.
\]

Из равенств
\[
AB=CD,
\qquad
BC=DE
\]
по свойствам равных хорд и дуг получаем
\[
AC=CE.
\]

Кроме того,
\[
AB=CD
\quad\Rightarrow\quad
\overset{\frown}{AB}
=
\overset{\frown}{CD}.
\]

Поэтому
\[
\angle ADB=\angle DBC.
\]

Следовательно,
\[
AD\parallel BC.
\]

Четырёхугольник \(ABCD\) является вписанной
равнобедренной трапецией, поэтому
\[
AC=BD.
\]

Применим теорему Птолемея к \(ABCD\):
\[
AC\cdot BD
=
AB\cdot CD
+
BC\cdot AD.
\]

Подставляем данные:
\[
AC\cdot BD
=
3\cdot3+4\cdot6
=
9+24
=
33.
\]

Так как
\[
AC=BD,
\]
получаем
\[
AC^2=33,
\]
следовательно,
\[
AC=BD=\sqrt{33}.
\]

Ранее установлено:
\[
CE=AC,
\]
поэтому
\[
CE=\sqrt{33}.
\]

Теперь рассмотрим вписанный четырёхугольник \(BCDE\).

По теореме Птолемея:
\[
BD\cdot CE
=
BC\cdot DE
+
CD\cdot BE.
\]

Подставляем:
\[
\sqrt{33}\cdot\sqrt{33}
=
4\cdot4+3\cdot BE.
\]

Получаем:
\[
33=16+3BE,
\]
\[
3BE=17,
\]
\[
\boxed{BE=\frac{17}{3}}.
\]

\section*{7. Рабочие алгоритмы}
\sectionrule

\subsection*{Если есть прямоугольный треугольник}

\begin{enumerate}

  \item Определите гипотенузу.

  \item Относительно нужного угла определите
  противолежащий и прилежащий катеты.

  \item Если стороны не хватает,
  примените теорему Пифагора.

  \item Выберите отношение:
  \[
  \sin,\qquad \cos,\qquad \tg.
  \]

  \item Сократите полученную дробь.

  \item Если формат экзамена требует десятичную запись,
  переведите результат в десятичную дробь в самом конце.

\end{enumerate}

\subsection*{Если есть касательная к окружности}

\begin{enumerate}

  \item Найдите точку касания.

  \item Проведите к ней радиус.

  \item Отметьте прямой угол:
  \[
  R\perp\text{касательной}.
  \]

  \item Ищите образовавшийся прямоугольный треугольник.

  \item Используйте центральные углы,
  диаметры и градусные меры дуг.

\end{enumerate}

\subsection*{Если есть равные хорды}

\begin{enumerate}

  \item Перейдите от равенства хорд
  к равенству соответствующих дуг.

  \item При необходимости сложите равные дуги.

  \item Перейдите обратно от равных дуг
  к равным хордам.

  \item Проверьте, возникают ли
  равные вписанные углы.

  \item Через равные углы
  проверьте возможность доказать параллельность прямых.

\end{enumerate}

\subsection*{Если есть вписанный четырёхугольник}

\begin{enumerate}

  \item Выпишите известные стороны и диагонали.

  \item Проверьте равенство хорд,
  дуг или диагоналей.

  \item Проверьте возможность применения
  теоремы Птолемея.

  \item Если в исходной фигуре можно выделить
  второй вписанный четырёхугольник,
  примените теорему Птолемея повторно.

\end{enumerate}

\section*{8. Типичные ошибки}
\sectionrule

\begin{itemize}

  \item Считать рисунок точным источником данных.
  Основной источник данных --- условие и доказанные свойства.

  \item Считать произвольный вписанный многоугольник
  правильным только из-за аккуратного чертежа.

  \item Путать противолежащий и прилежащий катеты
  в формулах синуса и косинуса.

  \item Писать «\(\sin\)», «\(\cos\)» или «\(\tg\)»
  без указания угла.

  \item Забывать, что радиус,
  проведённый в точку касания,
  перпендикулярен касательной.

  \item Путать центральный и вписанный углы:
  центральный угол равен дуге,
  вписанный угол равен половине дуги.

  \item Из равенства двух сторон
  сразу делать вывод о равенстве треугольников.

  \item Применять теорему Птолемея
  к четырёхугольнику,
  про который не установлено,
  что он вписан в окружность.

  \item Делать громоздкие вычисления,
  не проверив возможность вынести общий множитель,
  сгруппировать слагаемые
  или распознать полный квадрат.

  \item Использовать специальную формулу,
  не проверив её условия применимости.

\end{itemize}

\section*{9. Что особенно важно запомнить}
\sectionrule

\begin{enumerate}

  \item
  \[
  a^2+b^2=c^2
  \]
  для прямоугольного треугольника.

  \item
  \[
  \sin\alpha
  =
  \frac{\text{противолежащий катет}}
       {\text{гипотенуза}}.
  \]

  \item
  \[
  \cos\alpha
  =
  \frac{\text{прилежащий катет}}
       {\text{гипотенуза}}.
  \]

  \item Радиус к точке касания
  перпендикулярен касательной.

  \item Центральный угол равен дуге,
  на которую он опирается.

  \item Равные хорды
  стягивают равные дуги.

  \item Равные дуги
  стягиваются равными хордами.

  \item Вписанная трапеция
  является равнобедренной.

  \item Для вписанного четырёхугольника:
  \[
  AC\cdot BD
  =
  AB\cdot CD
  +
  BC\cdot AD.
  \]

\end{enumerate}

% ============================================================
% ТРЕНИРОВОЧНЫЙ БЛОК
% ============================================================

\newpage

\section*{Тренировочный блок}
\sectionrule

\textcolor{textgray}{
Решите самостоятельно.
Задания расположены по нарастающей сложности.
}

\subsection*{Средний уровень}

\begin{enumerate}[label=\textbf{\arabic*.}]

  \item
  В прямоугольном треугольнике \(ABC\)
  угол \(C\) прямой,
  \[
  AC=12,
  \qquad
  BC=5.
  \]
  Найдите
  \[
  \sin\angle BAC,
  \qquad
  \cos\angle BAC,
  \qquad
  \tg\angle BAC.
  \]

  \item
  В прямоугольном треугольнике \(MNP\)
  угол \(N\) прямой,
  \[
  MN=8,
  \qquad
  MP=17.
  \]
  Найдите \(NP\), а также
  \[
  \sin\angle NMP
  \quad\text{и}\quad
  \cos\angle NMP.
  \]

  \item
  Из точки \(C\) к окружности
  проведена касательная \(CA\).
  Прямая \(CO\) проходит через центр окружности
  и пересекает её в точках \(B\) и \(D\),
  причём точки расположены в порядке
  \[
  C,\ B,\ O,\ D.
  \]
  Дуга \(AD\) равна
  \[
  132^\circ.
  \]
  Найдите
  \[
  \angle ACO.
  \]

\end{enumerate}

\subsection*{Уровень выше среднего}

\begin{enumerate}[label=\textbf{\arabic*.},resume]

  \item
  В той же конфигурации
  дуга \(AD\) равна
  \[
  104^\circ.
  \]
  Найдите \(\angle ACO\) двумя способами:

  \begin{enumerate}[label=\arabic*)]
    \item через прямоугольный треугольник;
    \item через формулу угла между касательной и секущей.
  \end{enumerate}

  \item
  Пятиугольник \(ABCDE\)
  вписан в окружность.
  Известно:
  \[
  AB=CD,
  \qquad
  BC=DE.
  \]
  Докажите:
  \[
  AC=CE.
  \]

  \item
  Четырёхугольник \(ABCD\)
  вписан в окружность, причём
  \[
  AB=CD=5,
  \qquad
  BC=8,
  \qquad
  AD=12.
  \]
  Докажите:
  \[
  AD\parallel BC,
  \]
  а затем найдите
  \[
  AC
  \quad\text{и}\quad
  BD.
  \]

  \item
  Пятиугольник \(ABCDE\)
  вписан в окружность.
  Дано:
  \[
  AB=CD=3,
  \qquad
  BC=DE=5,
  \qquad
  AD=7.
  \]
  Найдите
  \[
  BE.
  \]

  \item
  Равнобедренная трапеция \(ABCD\)
  вписана в окружность,
  \[
  AD\parallel BC,
  \]
  причём
  \[
  AB=CD=6,
  \qquad
  AD=10,
  \qquad
  BC=4.
  \]
  Найдите обе диагонали.

  \item
  Из точки \(C\)
  к окружности проведена касательная \(CA\),
  а прямая \(CO\) проходит через центр
  и пересекает окружность в \(B\) и \(D\).

  Точки расположены в порядке
  \[
  C,\ B,\ O,\ D.
  \]

  Известно:
  \[
  \angle ACO=27^\circ.
  \]

  Найдите дугу
  \[
  \overset{\frown}{AD}.
  \]

\end{enumerate}

\subsection*{Сложный уровень}

\begin{enumerate}[label=\textbf{\arabic*.},resume]

  \item
  Пятиугольник \(ABCDE\)
  вписан в окружность.
  Известно:
  \[
  AB=CD=4,
  \qquad
  BC=DE=6,
  \qquad
  AD=8.
  \]

  \begin{enumerate}[label=\arabic*)]
    \item Докажите, что
    \[
    AC=CE.
    \]

    \item Найдите
    \[
    BE.
    \]
  \end{enumerate}

\end{enumerate}

% ============================================================
% ОТВЕТЫ
% ============================================================

\newpage

\section*{Ответы к тренировочному блоку}
\sectionrule

\renewcommand{\arraystretch}{1.35}

\begin{table}[h!]
\centering
\caption{Краткие ответы}

\begin{tabularx}{\linewidth}{@{}>{\bfseries}c X@{}}

\toprule
№ & Ответ \\
\midrule

1 &
\(
\sin\angle BAC=\dfrac{5}{13},
\quad
\cos\angle BAC=\dfrac{12}{13},
\quad
\tg\angle BAC=\dfrac{5}{12}
\).
\\

2 &
\(
NP=15,
\quad
\sin\angle NMP=\dfrac{15}{17},
\quad
\cos\angle NMP=\dfrac{8}{17}
\).
\\

3 &
\(
42^\circ
\).
\\

4 &
\(
14^\circ
\).
\\

5 &
\(
AC=CE
\).
\\

6 &
\(
AD\parallel BC,
\quad
AC=BD=11
\).
\\

7 &
\(
BE=\dfrac{19}{3}
\).
\\

8 &
\(
AC=BD=2\sqrt{19}
\).
\\

9 &
\(
\overset{\frown}{AD}=117^\circ
\).
\\

10 &
\(
AC=CE,
\quad
BE=7
\).
\\

\bottomrule

\end{tabularx}
\end{table}

\end{document}